%% kapa: exam=1 round=36 year=2025 period=2 form=A subject=appraisal_law subject_ko=감정평가관계법규 seq=4 number=4
%% source: Q-Net 감정평가사 기출문제 (cst00309) artlSeq=5231525 file=raw/1008/5231525/2025년 제36회 감정평가사 1차시험 2교시 문제지 원본.pdf
%% conversion: kordoc 4.24.51 -> tools/gen1.py (auto); review=proofread-sonnet; flags=none
\begin{question}{4}
국토의 계획 및 이용에 관한 법령상 공간재구조화계획에 관한 설명으로 옳지 않은 것은?
\begin{choices}
  \item 국토교통부장관은 특화발전 및 지역 균형발전을 위하여 필요한 때에는 관할 시장 또는 군수의 요청에 따라 공간재구조화계획을 입안할 수 있다.
  \item 주민이 복합용도구역의 지정을 위하여 공간재구조화계획의 입안을 제안하려면 대상 토지면적(국유지 포함)의 3분의 2 이상의 토지소유자의 동의를 받아야 한다.
  \item 주거지역에 공간재구조화계획을 입안하는 경우에 토지적성평가를 생략할 수 있다.
  \item 시ㆍ도지사가 결정하는 공간재구조화계획 중 복합용도구역 지정 및 입지 타당성 등에 관한 사항은 중앙도시계획위원회의 심의를 거친다.
  \item 지형도면이 필요 없는 경우에 공간재구조화계획 결정의 효력은 그 계획 결정을 고시한 날부터 발생한다.
\end{choices}
\end{question}
\answer{2}
